% Options for packages loaded elsewhere
\PassOptionsToPackage{unicode}{hyperref}
\PassOptionsToPackage{hyphens}{url}
\PassOptionsToPackage{dvipsnames,svgnames,x11names}{xcolor}
%
\documentclass[
  11pt,
  a4paper,
]{ltjsarticle}
\usepackage{amsmath,amssymb}
\usepackage{iftex}
\ifPDFTeX
  \usepackage[T1]{fontenc}
  \usepackage[utf8]{inputenc}
  \usepackage{textcomp} % provide euro and other symbols
\else % if luatex or xetex
  \usepackage{unicode-math} % this also loads fontspec
  \defaultfontfeatures{Scale=MatchLowercase}
  \defaultfontfeatures[\rmfamily]{Ligatures=TeX,Scale=1}
\fi
\usepackage{lmodern}
\ifPDFTeX\else
  % xetex/luatex font selection
\fi
% Use upquote if available, for straight quotes in verbatim environments
\IfFileExists{upquote.sty}{\usepackage{upquote}}{}
\IfFileExists{microtype.sty}{% use microtype if available
  \usepackage[]{microtype}
  \UseMicrotypeSet[protrusion]{basicmath} % disable protrusion for tt fonts
}{}
\makeatletter
\@ifundefined{KOMAClassName}{% if non-KOMA class
  \IfFileExists{parskip.sty}{%
    \usepackage{parskip}
  }{% else
    \setlength{\parindent}{0pt}
    \setlength{\parskip}{6pt plus 2pt minus 1pt}}
}{% if KOMA class
  \KOMAoptions{parskip=half}}
\makeatother
\usepackage{xcolor}
\usepackage[margin=24mm]{geometry}
\usepackage{color}
\usepackage{fancyvrb}
\newcommand{\VerbBar}{|}
\newcommand{\VERB}{\Verb[commandchars=\\\{\}]}
\DefineVerbatimEnvironment{Highlighting}{Verbatim}{commandchars=\\\{\}}
% Add ',fontsize=\small' for more characters per line
\newenvironment{Shaded}{}{}
\newcommand{\AlertTok}[1]{\textcolor[rgb]{1.00,0.00,0.00}{\textbf{#1}}}
\newcommand{\AnnotationTok}[1]{\textcolor[rgb]{0.38,0.63,0.69}{\textbf{\textit{#1}}}}
\newcommand{\AttributeTok}[1]{\textcolor[rgb]{0.49,0.56,0.16}{#1}}
\newcommand{\BaseNTok}[1]{\textcolor[rgb]{0.25,0.63,0.44}{#1}}
\newcommand{\BuiltInTok}[1]{\textcolor[rgb]{0.00,0.50,0.00}{#1}}
\newcommand{\CharTok}[1]{\textcolor[rgb]{0.25,0.44,0.63}{#1}}
\newcommand{\CommentTok}[1]{\textcolor[rgb]{0.38,0.63,0.69}{\textit{#1}}}
\newcommand{\CommentVarTok}[1]{\textcolor[rgb]{0.38,0.63,0.69}{\textbf{\textit{#1}}}}
\newcommand{\ConstantTok}[1]{\textcolor[rgb]{0.53,0.00,0.00}{#1}}
\newcommand{\ControlFlowTok}[1]{\textcolor[rgb]{0.00,0.44,0.13}{\textbf{#1}}}
\newcommand{\DataTypeTok}[1]{\textcolor[rgb]{0.56,0.13,0.00}{#1}}
\newcommand{\DecValTok}[1]{\textcolor[rgb]{0.25,0.63,0.44}{#1}}
\newcommand{\DocumentationTok}[1]{\textcolor[rgb]{0.73,0.13,0.13}{\textit{#1}}}
\newcommand{\ErrorTok}[1]{\textcolor[rgb]{1.00,0.00,0.00}{\textbf{#1}}}
\newcommand{\ExtensionTok}[1]{#1}
\newcommand{\FloatTok}[1]{\textcolor[rgb]{0.25,0.63,0.44}{#1}}
\newcommand{\FunctionTok}[1]{\textcolor[rgb]{0.02,0.16,0.49}{#1}}
\newcommand{\ImportTok}[1]{\textcolor[rgb]{0.00,0.50,0.00}{\textbf{#1}}}
\newcommand{\InformationTok}[1]{\textcolor[rgb]{0.38,0.63,0.69}{\textbf{\textit{#1}}}}
\newcommand{\KeywordTok}[1]{\textcolor[rgb]{0.00,0.44,0.13}{\textbf{#1}}}
\newcommand{\NormalTok}[1]{#1}
\newcommand{\OperatorTok}[1]{\textcolor[rgb]{0.40,0.40,0.40}{#1}}
\newcommand{\OtherTok}[1]{\textcolor[rgb]{0.00,0.44,0.13}{#1}}
\newcommand{\PreprocessorTok}[1]{\textcolor[rgb]{0.74,0.48,0.00}{#1}}
\newcommand{\RegionMarkerTok}[1]{#1}
\newcommand{\SpecialCharTok}[1]{\textcolor[rgb]{0.25,0.44,0.63}{#1}}
\newcommand{\SpecialStringTok}[1]{\textcolor[rgb]{0.73,0.40,0.53}{#1}}
\newcommand{\StringTok}[1]{\textcolor[rgb]{0.25,0.44,0.63}{#1}}
\newcommand{\VariableTok}[1]{\textcolor[rgb]{0.10,0.09,0.49}{#1}}
\newcommand{\VerbatimStringTok}[1]{\textcolor[rgb]{0.25,0.44,0.63}{#1}}
\newcommand{\WarningTok}[1]{\textcolor[rgb]{0.38,0.63,0.69}{\textbf{\textit{#1}}}}
\setlength{\emergencystretch}{3em} % prevent overfull lines
\providecommand{\tightlist}{%
  \setlength{\itemsep}{0pt}\setlength{\parskip}{0pt}}
\setcounter{secnumdepth}{5}
\ifLuaTeX
  \usepackage{selnolig}  % disable illegal ligatures
\fi
\IfFileExists{bookmark.sty}{\usepackage{bookmark}}{\usepackage{hyperref}}
\IfFileExists{xurl.sty}{\usepackage{xurl}}{} % add URL line breaks if available
\urlstyle{same}
\hypersetup{
  colorlinks=true,
  linkcolor={blue},
  filecolor={Maroon},
  citecolor={Blue},
  urlcolor={blue},
  pdfcreator={LaTeX via pandoc}}

\author{}
\date{}

\begin{document}

\hypertarget{ux7ddaux5f62ux4ee3ux6570ux30c6ux30adux30b9ux30c8}{%
\section{線形代数テキスト}\label{ux7ddaux5f62ux4ee3ux6570ux30c6ux30adux30b9ux30c8}}

この \texttt{docs/}
は、定義を順番に覚えるためではなく、\textbf{一つの概念が次の概念をなぜ必要とするか}を追うために構成されています。

\hypertarget{ux8aadux307fux65b9}{%
\subsection{読み方}\label{ux8aadux307fux65b9}}

各章では最初に \texttt{README.md} を読み、その後 numbered files
を順番に進めてください。

\begin{Shaded}
\begin{Highlighting}[]
\NormalTok{01 ベクトル空間}
\NormalTok{ ↓ 「空間をどう座標化するか」}
\NormalTok{02 線形写像}
\NormalTok{ ↓ 「写像の解をどう調べるか」}
\NormalTok{03 連立方程式}
\NormalTok{ ↓ 「近さ・直角をどう定義するか」}
\NormalTok{04 内積と幾何}
\NormalTok{ ↓}
\NormalTok{05 行列式 ── 可逆性を別の角度から理解}
\NormalTok{ ↓}
\NormalTok{06 固有値 ── 変換に自然な方向を探す}
\NormalTok{ ↓}
\NormalTok{07 行列分解 ── 複雑な写像を単純な部品に分ける}
\NormalTok{ ↓}
\NormalTok{08 最小二乗・逆問題}
\NormalTok{ ↓}
\NormalTok{09 低ランク・PCA}
\NormalTok{ ↓}
\NormalTok{10 発展理論 / 11 数値計算}
\end{Highlighting}
\end{Shaded}

\hypertarget{ux5404ux7bc0ux3067ux5fc5ux305aux78baux8a8dux3059ux308b4ux70b9}{%
\subsection{各節で必ず確認する4点}\label{ux5404ux7bc0ux3067ux5fc5ux305aux78baux8a8dux3059ux308b4ux70b9}}

\begin{enumerate}
\def\labelenumi{\arabic{enumi}.}
\tightlist
\item
  \textbf{定義} --- 数式として何を意味するか。
\item
  \textbf{幾何学} --- 空間の中で何が起きているか。
\item
  \textbf{代数} --- 他の定理・式からどう導かれるか。
\item
  \textbf{接続} --- 後のどの概念に必要になるか。
\end{enumerate}

例えば SVD を学ぶとき、\texttt{A=UΣV\^{}T}
を暗記するだけでは不十分です。\texttt{V}
がなぜ入力側の正規直交基底なのか、\texttt{σ\_i}
がなぜ伸縮率なのか、零特異値がなぜ核を表すのか、そこから擬似逆や PCA
がなぜ出るのかを説明できる状態を目標にします。

\hypertarget{ux7406ux89e3ux306eux5224ux5b9aux57faux6e96}{%
\subsection{理解の判定基準}\label{ux7406ux89e3ux306eux5224ux5b9aux57faux6e96}}

ある概念について次の3問に答えられれば、かなり深く理解できています。

\begin{itemize}
\tightlist
\item
  その定義を使わず、直感的な日本語で説明できるか。
\item
  その定義が必要になる具体例を一つ作れるか。
\item
  一つ前と一つ後の概念との関係を説明できるか。
\end{itemize}

計算問題を解けることは必要ですが、それだけを到達点にはしません。

\end{document}
